\writeSectionFrame{\presentationTitle}{Fazit}
\begin{frame}[fragile]{ArchUnit als Beispiel}
	\begin{lstlisting}[basicstyle=\tiny]
ArchRule domain_is_framework_independent =
    noClasses()
        .that().resideInAPackage("..domain..")
        .should().dependOnClassesThat()
        .resideInAnyPackage(
            "org.springframework..",
            "jakarta..",
            "javax.."
        );

	\end{lstlisting}
\end{frame}

\begin{frame}[fragile]{ArchUnit als Beispiel}
	\begin{lstlisting}[basicstyle=\tiny]
JavaClasses classes =
    new ClassFileImporter().importPackages("de.example.app");

ArchRule rule =
    layeredArchitecture()
        .layer("Controller").definedBy("..controller..")
        .layer("Service").definedBy("..service..")
        .layer("Domain").definedBy("..domain..")
        .layer("Persistence").definedBy("..persistence..")

        .whereLayer("Controller").mayOnlyBeAccessedByLayers("Service")
        .whereLayer("Service").mayOnlyBeAccessedByLayers("Controller")
        .whereLayer("Domain").mayOnlyBeAccessedByLayers("Service", "Persistence")
        .whereLayer("Persistence").mayNotBeAccessedByAnyLayer();

rule.check(classes);
	\end{lstlisting}
\end{frame}

\begin{frame}{Vorteile}
	\begin{itemize}
		\item Sehr gute Lesbarkeit
		\item Weniger temporäre Variablen
		\item Führt den Nutzer durch erlaubte Aufrufreihenfolgen
	\end{itemize}
\end{frame}


\begin{frame}{Nachteile und Risiken}
	\begin{itemize}
		\item Kann komplexe Logik verschleiern
		\item Debugging manchmal schwieriger
		\item Gefahr von \glqq API-Spielereien\grqq
		\item Gute Namenswahl entscheidend
	\end{itemize}
\end{frame}

\begin{frame}{Regel}
\centering
\Large
\emph{Eine Fluent API ist dann gut, wenn der Code\\
seine eigene Dokumentation ist.}

\emph{Eine gute Fluent API zwingt uns nicht, Code zu lesen –
sie erlaubt uns, Absichten zu lesen.}
\end{frame}
